% Generated draft: requires mathematical and source review.
\documentclass[11pt]{article}
\usepackage[margin=25mm]{geometry}
\usepackage{fontspec}
\setmainfont{Latin Modern Roman}
\usepackage{amsmath,amssymb,mathtools,unicode-math}
\setmathfont{Latin Modern Math}
\usepackage{enumitem,needspace,xurl,hyperref}
\setlength{\parindent}{0pt}
\setlength{\parskip}{5pt}
\setlength{\emergencystretch}{3em}
\newcommand{\N}{\mathbb{N}}
\newcommand{\Z}{\mathbb{Z}}
\newcommand{\Q}{\mathbb{Q}}
\newcommand{\R}{\mathbb{R}}
\newcommand{\C}{\mathbb{C}}
\newcommand{\F}{\mathbb{F}}
\newcommand{\A}{\mathbb{A}}
\newcommand{\PP}{\mathbb{P}}
\newcommand{\E}{\mathbb{E}}
\newcommand{\eps}{\varepsilon}
\newcommand{\dd}{\,\mathrm{d}}
\begin{document}
\section*{UnsolvedMath ID: 1101223}
\textbf{Source problem code:} catalogue entry
\subsection{Definitions and mathematical statement}
Fix an integer $n\geq1$. The braid group $B_n$ has generators $\sigma_1,\ldots,\sigma_{n-1}$, with commuting relations for indices at distance greater than one and the braid relations $\sigma_i\sigma_{i+1}\sigma_i=\sigma_{i+1}\sigma_i\sigma_{i+1}$. The singular braid monoid $SB_n$ has generators $\sigma_i^{\pm1},\tau_i$ for $1\leq i<n$. In addition to the braid and inverse relations for the $\sigma_i$, its relations are
\[
\begin{aligned}
 \tau_i\tau_j&=\tau_j\tau_i &&(|i-j|>1),\\
 \sigma_i\tau_j&=\tau_j\sigma_i &&(|i-j|>1),\\
 \sigma_i\tau_i&=\tau_i\sigma_i,\\
 \sigma_i\sigma_{i+1}\tau_i&=\tau_{i+1}\sigma_i\sigma_{i+1} &&(1\leq i<n-1),\\
 \sigma_{i+1}\sigma_i\tau_{i+1}&=\tau_i\sigma_{i+1}\sigma_i &&(1\leq i<n-1).
\end{aligned}
\]
Here $\tau_i$ represents a singular crossing between strands $i$ and $i+1$. The integral group ring $\Z[B_n]$ consists of finite formal integer linear combinations of elements of $B_n$, with multiplication extending the group multiplication bilinearly.

The desingularization map is the monoid homomorphism to the multiplicative monoid of $\Z[B_n]$ defined by
\[
 \Phi:SB_n\longrightarrow\Z[B_n],\qquad
 \Phi(\sigma_i^{\pm1})=\sigma_i^{\pm1},\qquad
 \Phi(\tau_i)=\sigma_i-\sigma_i^{-1}.
\]
\textbf{Recorded conjecture.} For every $n\geq1$ and every $x,y\in SB_n$,
\[
 \Phi(x)=\Phi(y)\ \Longrightarrow\ x=y.
\]
The supplied research summary uses the symbol $\eta$ for this same map and reports its injectivity for all $n$, attributing the result to Paris (2004). It further reports the consequence that Vassiliev braid invariants classify singular braids. These are supplied solved-status claims, not independent verification in this draft.
\subsection{Short English statement}
The desingularization map replaces each singular crossing by the difference between its positive and negative resolutions. The question asks whether its value in the integral braid-group ring determines the singular braid uniquely, for every number of strands. The record attributes a positive answer to Paris (2004).
\subsection{Sources}
\begin{itemize}
\item Local UnsolvedMath registry, permanent ID 1101223, source code catalogue entry, ``Questions in Geometric Group Theory --- Q 12.23''; supplied status solved, classification SOLVED-IN-LITERATURE, updated 2026-08-06.
\item Mladen Bestvina, \emph{Questions in Geometric Group Theory}, list dated 2004, Question 12.23, PDF page 24, \url{https://www.math.utah.edu/~bestvina/eprints/questions-updated.pdf}. Reported access: 2026-07-29; extraction: manual transcription. The background says that the author PDF records partial progress or special cases and current status still requires release review. Rights: NEEDS\_REVIEW, public author-hosted PDF.
\item The supplied research summary attributes injectivity of the desingularization map to Paris (2004); no title, journal, theorem number, or other bibliographic detail is supplied. It reports the consequence concerning Vassiliev braid invariants.
\end{itemize}
\end{document}
